%%%%%%%%%%%%%%%%%%%%%%%%%%%%%%%%%%%%%%%%%%%%%%%%%%%%%%%%%%%%%%%%%%%%%%%%
% Refine As You Round — MLSys 2027 submission (research track)
% Format: mlsys2025style (MLSys 2027 uses the 2025 style; 2-column,
% <= 10 pages excluding references; appendix uploaded separately).
%
% Build:  latexmk -pdf main.tex
% Anonymous submission: \usepackage{mlsys2025}
% Camera-ready:         \usepackage[accepted]{mlsys2025}
%%%%%%%%%%%%%%%%%%%%%%%%%%%%%%%%%%%%%%%%%%%%%%%%%%%%%%%%%%%%%%%%%%%%%%%%
\documentclass{article}

\usepackage{microtype}
\usepackage{graphicx}
\usepackage{subfigure}
\usepackage{booktabs}
\usepackage{multirow}
\usepackage{amsmath,amssymb}
\usepackage{xcolor}
\usepackage{hyperref}

% Anonymous submission (double-blind). Switch to [accepted] for camera-ready.
\usepackage{mlsys2025}

% ---- placeholder macros: red = numbers to fill, blue = systems co-author
\newcommand{\tbd}[1]{\textcolor{red}{[#1]}}
\newcommand{\teammate}[1]{\textcolor{blue}{[TEAMMATE: #1]}}
\newcommand{\blank}{\textcolor{red}{[\;\;]}}

\newcommand{\bpw}{\mathrm{bpw}}
\newcommand{\gain}{\mathrm{gain}}
\newcommand{\tr}{\mathrm{tr}}
\newcommand{\Hin}{H_{\mathrm{in}}}
\newcommand{\Hout}{H_{\mathrm{out}}}

\mlsystitlerunning{Refine As You Round: In-Loop Nested-Lattice Mixed Precision}

\begin{document}

\twocolumn[
\mlsystitle{Refine As You Round: In-Loop Nested-Lattice Mixed Precision\\
for Entropy-Coded 3-bit LLM Weights}

\mlsyssetsymbol{equal}{*}

\begin{mlsysauthorlist}
\mlsysauthor{Anonymous Authors}{anon}
\end{mlsysauthorlist}

\mlsysaffiliation{anon}{Anonymous Institution, Anonymous City, Anonymous Region, Anonymous Country}

\mlsyscorrespondingauthor{Anonymous Author}{anon.email@domain.com}

\mlsyskeywords{Machine Learning, MLSys, quantization, mixed precision, entropy coding}

\vskip 0.3in

\begin{abstract}
Weight-only post-training quantization (PTQ) of large language models at
3 bits per weight still loses noticeably more perplexity than 4 bits,
while 4 bits wastes rate on the majority of weights that are already well
represented at 3. Mixed precision at element granularity closes this gap
in principle, but existing schemes either store an explicit sparse index
for the high-precision elements (SpQR) or need a full gradient pass over
the model to pick them (SqueezeLLM). We present a representation and a
selection rule that need neither. The representation is a nested lattice:
the 3-bit grid of every quantization group is embedded in a 4-bit grid
with half the step and the same zero, so the weight stream is a single
4-bit code stream in which refined elements are the odd codes; the mask
is implicit and the stream has low entropy, so a per-tensor Huffman code
stores it at $\approx 3.07$ bits/weight including all side information.
The selection is made inside the GPTQ-style column loop, at the moment
each weight is rounded, from the solver's own second-order error weight
applied to the error-compensated weight: $\gain_{ij} = (e_c^2 - e_f^2)\,
d_j\, r_i$, where $d_j$ and $r_i$ are the column and row weights the
solver already uses (two-sided where the solver is two-sided). One global
threshold on the layer-normalised gain allocates the refinement budget
across the whole model. On Llama-3.2-1B quantised with TurboBoA at
$g{=}128$ and matched entropy-coded rate, the in-loop rule
\tbd{matches/beats} the static diagonal-Fisher mask by \tbd{$\Delta$
wiki2 PPL over 3 seeds}, with zero extra passes over the model; the rule
transfers to \tbd{3B/8B and to Mistral/Qwen} \tbd{one clause}. The second
half of the paper describes the bitstream format, the decoder, and the
fused dequantisation kernels that make the entropy-coded stream usable at
inference time \teammate{one-sentence systems result}.
\end{abstract}
]

\printAffiliationsAndNotice{}

%%%%%%%%%%%%%%%%%%%%%%%%%%%%%%%%%%%%%%%%%%%%%%%%%%%%%%%%%%%%%%%%%%%%%%%%
\section{Introduction}
\label{sec:intro}

Serving a large language model (LLM) at batch size one is bound by the
bytes of weights that must cross the memory bus per generated token, so
the effective bits per weight (bpw) of the stored model translates almost
linearly into decode latency and into how large a model fits on a given
device. Post-training quantization (PTQ) of weights to 4 bits with
per-group scales is now routine and nearly lossless for models above a
few billion parameters \citep{frantar2023gptq,lin2024awq}; 3 bits is
where the loss becomes visible, and 2 bits is where it becomes severe
\citep{chee2023quip,egiazarian2024aqlm}. Between 3 and 4 bits sits a
well-known asymmetry: only a small fraction of the weights are
responsible for most of the 3-bit degradation, and giving those weights
one more bit recovers a disproportionate share of the 4-bit quality
\citep{dettmers2024spqr,kim2024squeezellm}.

Exploiting that asymmetry at element granularity raises two questions
that existing methods answer expensively. \emph{Where is the extra
precision stored?} SpQR \citep{dettmers2024spqr} keeps the sensitive
weights in a separate sparse FP16 structure with explicit row/column
indices, which costs $\sim$32 bits per outlier and a second data path in
the kernel. \emph{How are the sensitive weights chosen?} SqueezeLLM
\citep{kim2024squeezellm} ranks weights by a diagonal empirical Fisher,
which needs a full backward pass over the model on the calibration set
before quantization starts; SpQR uses the quantizer's own error but
stores it densely.

This paper answers both questions with one construction. First, the
\emph{nested lattice}: within each quantization group the 3-bit grid
$\{s_c(q - z_c)\}$ is a sub-lattice of the 4-bit grid with step $s_c/2$
and zero $2z_c$. Every weight is stored as a 4-bit code; a 3-bit weight
is the even code $2q_c$ and a refined weight is any code $q_f$.
Dequantization is one multiply--add for every element, there is no mask,
no index and no second scale, and the code stream has low entropy (half
the codes are even and the histogram is peaked), so a 16-symbol Huffman
table per tensor brings the stored rate to $3+\epsilon$ bpw where
$\epsilon$ is set by the refined fraction. Second, the \emph{in-loop
selection rule}: the GPTQ/OBS family of solvers
\citep{frantar2022obc,frantar2023gptq} quantizes one column at a time on
the error-compensated weight $w$ and turns the rounding error into a loss
increment $(w-q)^2 d_j$ with a column weight $d_j$ derived from the
Cholesky factor of the inverse input Hessian. At the moment column $j$ is
rounded we therefore know, for every row $i$, exactly how much loss the
solver would save by rounding to the fine grid instead: $\gain_{ij} =
(e_c^2 - e_f^2)\, d_j\, r_i$, with $r_i$ the row weight in solvers that
carry an output-side Hessian. We refine iff the gain is large --- top-$k$
per column, or above one global threshold $\lambda$ on the gain
normalised by the layer's own loss scale --- and propagate the error of
the value that is actually stored. No gradients, no extra calibration
pass, no change to the solver outside the rounding call.

The natural baseline is the same nested representation with the mask
chosen by the static Fisher score. We run the comparison at matched
entropy-coded rate on Llama-3.2-1B \citep{grattafiori2024llama3} with
TurboBoA \citep{kim2025boa}, a two-sided GPTQ-family solver, across three
calibration seeds, and \tbd{state the outcome}. A surprising diagnostic
explains why the in-loop rule can win: \tbd{X}\% of the elements a
Fisher mask selects are already exactly representable on the coarse grid
(their fine code is even), so the static mask spends \tbd{half} of its
budget on elements where refinement changes nothing, whereas the in-loop
rule sees the compensated residual and never does. We then show that the
same rule and the same protocol transfer across model sizes
(Section~\ref{sec:scaling}), model families (Section~\ref{sec:families})
and against published 3-bit methods on the standard Llama-2 test bed
(Section~\ref{sec:published}).

\paragraph{Contributions.}
\begin{itemize}\itemsep 0pt
\item A nested-lattice mixed-precision representation for grouped weight
  quantization in which the per-element precision mask is implicit in a
  single entropy-coded code stream (Section~\ref{sec:representation}).
\item An in-loop refinement rule that reuses the solver's second-order
  error weights, including the row weight of two-sided (attention-aware)
  solvers, and a rate-only calibration of its single global threshold
  (Section~\ref{sec:inloop}).
\item An evaluation protocol that compares selection rules at matched
  entropy-coded bits/weight rather than nominal bits, with endpoints that
  verify the modified solver is a bit-exact superset of the stock one
  (Section~\ref{sec:protocol}), and results on Llama-3.2-1B at
  3$\to$4 and 2$\to$3 bits, on 3B/8B models, on other families, and
  against published methods (Section~\ref{sec:results}).
\item \teammate{The bitstream format, Huffman decoder and fused
  dequantisation kernel, with end-to-end decode throughput and memory
  numbers (Sections~\ref{sec:bitstream}--\ref{sec:sysval}).}
\end{itemize}

%%%%%%%%%%%%%%%%%%%%%%%%%%%%%%%%%%%%%%%%%%%%%%%%%%%%%%%%%%%%%%%%%%%%%%%%
\section{Background}
\label{sec:background}

\subsection{Weight-only PTQ with second-order error compensation}

For a linear layer with weight $W \in \mathbb{R}^{d_{\mathrm{out}}\times
d_{\mathrm{in}}}$ and calibration inputs $X$, GPTQ-family methods
minimise the layer output error $\mathcal{L}(\hat W) = \tr\big((\hat W -
W)\,\Hin\,(\hat W - W)^\top\big)$, $\Hin = XX^\top$, subject to $\hat W$
lying on the quantization grid. Optimal Brain Surgeon
\citep{hassibi1993obs} gives the exact update for quantizing one weight
and compensating the rest; GPTQ \citep{frantar2023gptq} applies it column
by column using the upper Cholesky factor $U$ of $\Hin^{-1}$: after
quantizing column $j$ of the compensated weight to $q$, the error
$\mathrm{err} = (w-q)/U_{jj}$ is propagated to the remaining columns as
$\mathrm{err}\,U_{j,j:}$, and the loss increases by exactly
$\mathrm{err}^2 = (w-q)^2/U_{jj}^2$. We write $d_j = 1/U_{jj}^2$ for
this column weight.

\subsection{What the solver minimises, layer by layer}
\label{sec:objectives}

The in-loop rule reuses the solver's loss weights, so we state exactly
what TurboBoA \citep{kim2025boa} minimises for each linear layer of a
Llama block; the per-element gain of Section~\ref{sec:inloop} is the
increment of these objectives. Notation: $X \in \mathbb{R}^{L\times d}$
holds the $L$ calibration tokens entering the layer as rows, $W \in
\mathbb{R}^{d_{\mathrm{out}}\times d_{\mathrm{in}}}$, $\Delta W = W -
\hat W$, and every covariance is an average over the calibration set (the
normalisation cancels in the solver). Llama-3.2-1B has $n_h = 32$ query
heads, $n_{kv} = 8$ key/value heads (each shared by $n_h/n_{kv} = 4$
query heads), head dimension $d_h = 64$, $d = 2048$ and MLP width 8192.
$W_{\cdot,h}$ denotes the $d_h$ rows of a projection that belong to head
$h$.

\paragraph{One-sided layers (o\_proj, gate, up, down).} The layer output
error is $\|X\Delta W^\top\|_F^2$, so
\begin{equation}
  \mathcal{L}_{1s}(\Delta W) = \tr\big(\Delta W\,\Hin\,\Delta W^\top\big),
  \qquad \Hin = X^\top X,
  \label{eq:onesided}
\end{equation}
with $X$ the concatenated head outputs for o\_proj, the post-attention
RMSNorm output for gate and up (which share one $\Hin$), and
$\mathrm{SiLU}(XW_{\mathrm{gate}}^\top)\odot XW_{\mathrm{up}}^\top$ for
down. This is the GPTQ objective.

\paragraph{Value projection (v\_proj; one-sided, per-head Hessian).}
Head $h$ of the attention output is $A_h X W_{v,h}^\top$ with $A_h \in
\mathbb{R}^{L\times L}$ the (causal) softmax attention matrix of the
quantised model, so the error that reaches the next layer is $A_h X\Delta
W_{v,h}^\top$ and
\begin{equation}
\begin{aligned}
  \mathcal{L}_v(\Delta W_v) &= \sum_{h=1}^{n_{kv}}
  \tr\big(\Delta W_{v,h}\,\Hin^{(h)}\,\Delta W_{v,h}^\top\big),\\
  \Hin^{(h)} &= X^\top A_h^\top A_h X,
\end{aligned}
\label{eq:value}
\end{equation}
where, under grouped-query attention, $A_h$ of a key/value head is the
average over the four query heads that share it (\texttt{block\_v} in
the solver).

\paragraph{Query and key projections (q\_proj, k\_proj; two-sided).}
Only the attention scores depend on $W_q$ and $W_k$. With RoPE, the score
between positions $t$ and $s$ in head $h$ is $s_{ts} = (R_t W_{q,h}
x_t)^\top (R_s W_{k,h} x_s)/\sqrt{d_h}$, where $R_t$ is the $d_h\times
d_h$ rotation of position $t$. Perturbing $W_{q,h}$ changes every score
in row $t$ by $\Delta s_{ts} = (R_t\Delta W_{q,h} x_t)^\top \tilde
k_{h,s}$, with $\tilde k_{h,s} = R_s W_{k,h} x_s$ the post-RoPE key, and
\begin{equation}
\begin{aligned}
  \sum_s \Delta s_{ts}^2 &= x_t^\top \Delta W_{q,h}^\top R_t^\top C_{K,h}
  R_t \Delta W_{q,h} x_t,\\
  C_{K,h} &= \sum_s \tilde k_{h,s}\tilde k_{h,s}^\top .
\end{aligned}
\label{eq:scores}
\end{equation}
Summing over $t$ and decoupling the position-dependent rotation from the
input covariance (the BoA approximation)
\begin{equation}
\begin{aligned}
  &\sum_t \tr\big(R_t^\top C R_t\,\Delta W x_t x_t^\top \Delta W^\top\big)\\
  &\qquad\approx \tr\Big(\big(\tfrac{1}{L}\textstyle\sum_t R_t^\top C R_t\big)
  \Delta W X^\top X \Delta W^\top\Big)
\end{aligned}
\label{eq:boaapprox}
\end{equation}
gives the two-sided objective
\begin{align}
  \mathcal{L}_q(\Delta W_q) &= \sum_{h=1}^{n_h}
  \tr\big(\Hout^{(h)}\Delta W_{q,h}\Hin\Delta W_{q,h}^\top\big),\nonumber\\
  \Hout^{(h)} &= \tfrac{1}{L}\textstyle\sum_{t=1}^{L} R_t^\top
  C_{K,g(h)} R_t, \label{eq:lq}\\
  \mathcal{L}_k(\Delta W_k) &= \sum_{h=1}^{n_{kv}}
  \tr\big(\Hout^{(h)}\Delta W_{k,h}\Hin\Delta W_{k,h}^\top\big),\nonumber\\
  \Hout^{(h)} &= \tfrac{1}{L}\textstyle\sum_{t=1}^{L} R_t^\top
  \bar C_{Q,h} R_t, \label{eq:lk}
\end{align}
where $\Hin = X^\top X$ is the shared input covariance of q/k/v, $g(h)$
is the key head serving query head $h$, and $\bar C_{Q,h}$ is the
post-RoPE query covariance averaged over the four query heads that share
key head $h$. The output-side factor is what makes these layers
two-sided: an error in row $i$ of head $h$ is weighted by $\Hout^{(h)}$,
not by the identity. The solver handles it with a row-side Cholesky
factor $U_{\mathrm{out}}$ of $(\Hout^{(h)})^{-1}$ (after 1\% damping) and
quantises rows in chunks of $n_r = 16$, applying the GPTQ column loop to
each chunk and propagating the chunk's error to the remaining rows; the
OBS weight of a chunk is $(U_{\mathrm{sub}}^\top U_{\mathrm{sub}})^{-1}$,
whose diagonal is our row weight $r_i$ (Section~\ref{sec:inloop}).

\paragraph{Input-difference correction.} Because earlier blocks are
already quantised, the layer sees $\hat X = X - \Delta X$ at inference.
TurboBoA minimises $\|XW^\top - \hat X\hat W^\top\|_F^2 = \tr(\Delta W
\hat X^\top\hat X\Delta W^\top) + 2\tr(\Delta W\hat X^\top\Delta X
W^\top) + \mathrm{const}$ with all Hessians above built from $\hat X$ and
the cross term damped by $\alpha = 0.25$ (\texttt{consider\_dX}). The
cross term is linear in $\Delta W$; in the column loop it shifts the
compensated target of the not-yet-quantised columns (the $-wP_{j,j:}$
term of the update) but does not depend on which grid level column $j$ is
rounded to, so the loss increment of the rounding decision itself is the
quadratic term alone: $(w_{ij} - \hat w_{ij})^2 d_j r_i$. That is the
quantity our gain compares between the coarse and the fine level, and
$N_\ell$ is the same quadratic form evaluated on $W$.

\paragraph{Scale search.} Per group the solver picks $(s_c, z_c)$ by a
grid search over shrink factors of the min--max range that minimises the
Hessian-weighted error $\tr(\Delta W_g \Hin^{gg}\Delta W_g^\top)$ of the
group's columns (\texttt{qparam\_comput=Hessian}); for the two-sided
layers it repeats the search on each row chunk after the row-side
compensation (\texttt{adaptive\_qparam}). We keep this search exactly as
shipped and restrict it to the coarse grid.

\subsection{Grouped quantization and side information}

With group size $g$, each run of $g$ input columns in a row has a scale
$s$ (16 bit) and an asymmetric zero $z$ ($b$ bits for a $b$-bit grid). At
$g{=}128$ this adds $(16+b)/128 \approx 0.15$ bpw, so ``3-bit'' weights
are 3.15 bpw nominally. All rates in this paper include it.

\subsection{Entropy coding of code streams}

Quantized codes are not uniformly distributed; Deep Compression
\citep{han2016deepcompression} already Huffman-coded them. With a
per-tensor table of $2^b$ code lengths the overhead is $4\cdot 2^b$ bits
per tensor, negligible against $10^7$ codes. We report the ideal
(Shannon) rate and the realised Huffman rate \citep{huffman1952}; an ANS
coder \citep{duda2013ans} would approach the former, and
Section~\ref{sec:bitstream} describes the decoder we actually built.

\subsection{Nested lattices}

A lattice $\Lambda_c$ is nested in $\Lambda_f$ if $\Lambda_c \subset
\Lambda_f$ \citep{conway1982lattice}. For scalar uniform grids this is
just ``the coarse step is a multiple of the fine step and the zeros
coincide''. NestQuant \citep{savkin2025nestquant} uses nested $E_8$
lattices for matrix-product quantization with one fixed lattice pair per
tensor; Matryoshka-style methods
\citep{nair2025matquant,kleinegger2026matgptq,park2024anyprecision} nest
bit-widths by MSB slicing of a single master code so that one stored model
can be served at several uniform precisions. We use nesting differently:
to give individual elements one extra bit inside one code stream.

%%%%%%%%%%%%%%%%%%%%%%%%%%%%%%%%%%%%%%%%%%%%%%%%%%%%%%%%%%%%%%%%%%%%%%%%
\section{The nested-lattice representation}
\label{sec:representation}

\paragraph{Grids.} Fix a quantization group and let $(s_c, z_c)$ be the
scale and zero found by the solver's own search on the $b$-bit grid ($b =
3$ in the main experiment, $z_c \in \{0,\dots,2^b-1\}$). Define the
coarse and fine quantizers
\begin{equation}
\begin{aligned}
  q_c &= \mathrm{clamp}\big(\mathrm{round}(w/s_c) + z_c,\, 0,\, 2^b{-}1\big),\\
  \mathrm{deq}_c &= s_c(q_c - z_c),\\
  s_f &= s_c/2,\qquad z_f = 2z_c,\\
  q_f &= \mathrm{clamp}\big(\mathrm{round}(w/s_f) + z_f,\, 0,\, 2^{b+1}{-}1\big),\\
  \mathrm{deq}_f &= s_f(q_f - z_f).
\end{aligned}
\label{eq:grids}
\end{equation}
Every coarse level is a fine level: $s_c(q_c - z_c) = s_f(2q_c - z_f)$
exactly, also in floating point, since halving a scale and doubling an
integer are exact.

\paragraph{Stored code.} Every element stores one $(b{+}1)$-bit integer
$q \in \{0,\dots,2^{b+1}-1\}$: $q = 2q_c$ if the element is coarse, $q =
q_f$ if it is refined. Dequantization is the same for all elements, $\hat
w = s_f(q - z_f)$, so a kernel never branches on the mask; one scale and
one $b$-bit zero per group are the only side information. The mask is
implicit --- refined elements are those with odd $q$ or an even $q$ that
happens to equal $2q_c$ (Section~\ref{sec:free}).

\paragraph{Scale search stays on the coarse grid.} $(s_c, z_c)$ are
searched on the $b$-bit grid on the original group, before any column of
it is quantised, and $(s_f, z_f)$ are derived. Re-searching on the fine
grid would break the nesting (the coarse levels would no longer be a
subset), and we show in Section~\ref{sec:results} that the derived fine
grid costs \tbd{+0.31 PPL (e4\_s0 vs.\ w4\_s0)} relative to a native
$(b{+}1)$-bit search when every element is refined --- the price of the
nesting, paid only by the refined 1\%.

\paragraph{Rate.} Let $p_{\mathrm{odd}}$ be the fraction of odd codes.
The code stream's entropy is at most $H(\text{coarse histogram}) +
h_2(p_{\mathrm{odd}}) + p_{\mathrm{odd}}$ bits, i.e.\ the coarse cost plus
a small term that vanishes as $p_{\mathrm{odd}}\to 0$; at
$p_{\mathrm{odd}} = 1\%$ it is $\approx 0.09$ bpw. Per tensor we build
one Huffman table over the $2^{b+1}$ symbols and add table bits, scales
and zeros:
\begin{equation}
  \bpw_{\mathrm{Huff}} = \frac{\sum_\ell\Big(\sum_q n_\ell(q) L_\ell(q) +
  2^{b+3} + G_\ell (16+b)\Big)}{\sum_\ell N_\ell},
  \label{eq:rate}
\end{equation}
where $n_\ell(q)$ is the count of code $q$ in tensor $\ell$, $L_\ell(q)$
its Huffman length, $G_\ell$ the number of groups and $N_\ell$ the
number of elements. All comparisons in this paper are made at matched
$\bpw_{\mathrm{Huff}}$; the Shannon rate is reported alongside.

\subsection{Free elements}
\label{sec:free}

For a coarse element, refinement changes the stored value only if $q_f
\neq 2q_c$, i.e.\ if the compensated weight lies closer to a fine level
that is not a coarse level. For a weight uniformly distributed inside a
coarse cell this happens with probability $1/2$. On Llama-3.2-1B we
measure \tbd{51.4\% (e3\_s0)} of all elements to be ``free'': their
fine code already equals $2q_c$, and refining them is a no-op that costs
nothing and gains nothing. A selection rule that does not look at the
residual --- any static sensitivity score --- cannot see this, and will
spend roughly half its budget on free elements. This single observation
is, we believe, the main reason the in-loop rule is competitive
(Section~\ref{sec:results}).

\begin{algorithm}[t]
\caption{In-loop nested refinement inside a GPTQ-style column loop (one
layer)}
\label{alg:inloop}
\begin{algorithmic}[1]
\STATE build $\Hin$ (and $\Hout$ where the solver does); $U \leftarrow
  \mathrm{chol}(\Hin^{-1})$; $d_j \leftarrow 1/U_{jj}^2$
\STATE per group: search $(s_c, z_c)$ on the $b$-bit grid; $s_f
  \leftarrow s_c/2$, $z_f \leftarrow 2z_c$
\STATE $N_\ell \leftarrow \mathcal{L}(W)/\mathrm{numel}(W)$ (global-$\lambda$
  rule only)
\FOR{each row chunk (two-sided) or the whole tensor (one-sided)}
  \STATE $r_i \leftarrow [(U_{\mathrm{sub}}^\top U_{\mathrm{sub}})^{-1}]_{ii}$
    (two-sided) or $1$
  \FOR{column $j$ in solver order, on the compensated $W$}
    \STATE $w \leftarrow W_{:,j}$; compute $q_c, q_f$, $e_c = w -
      \mathrm{deq}_c$, $e_f = w - \mathrm{deq}_f$
    \STATE $\gain \leftarrow \max\big((e_c^2 - e_f^2)\, d_j\, r_i,\, 0\big)$
      (vector over rows)
    \STATE $\mathrm{refine} \leftarrow \textsc{Select}(\gain)$
    \STATE $q \leftarrow \mathrm{refine}\ ?\ q_f : 2q_c$; $\hat w
      \leftarrow s_f(q - z_f)$
    \STATE $\mathrm{err} \leftarrow (w - \hat w)/U_{jj}$; propagate err
      exactly as the solver does
  \ENDFOR
\ENDFOR
\end{algorithmic}
\end{algorithm}

%%%%%%%%%%%%%%%%%%%%%%%%%%%%%%%%%%%%%%%%%%%%%%%%%%%%%%%%%%%%%%%%%%%%%%%%
\section{In-loop refinement}
\label{sec:inloop}

Only the rounding call inside the column loop changes; the Hessian
construction, Cholesky factorisation, block updates, the solver's
input-difference correction and its per-group scale search are exactly
as shipped (Algorithm~\ref{alg:inloop}).

\paragraph{The gain.} At the moment column $j$ is rounded, the solver's
loss increment for storing $\hat w$ in row $i$ is $(w_{ij} - \hat
w_{ij})^2 d_j r_i$ under each of the objectives
\eqref{eq:onesided}--\eqref{eq:lk} (Section~\ref{sec:objectives}). The
gain of refining is the difference between the coarse and fine
increments,
\begin{equation}
  \gain_{ij} = (e_c^2 - e_f^2)\, d_j\, r_i \;\ge\; 0,
  \label{eq:gain}
\end{equation}
which is in the solver's own loss units and accounts for the error
compensation already applied to $w$ by previous columns. Because
$\Lambda_c \subset \Lambda_f$, $e_f^2 \le e_c^2$ except for clamp edge
cases, which we clip to zero. The gain is zero exactly for free elements.

\paragraph{Row weight $r_i$.} For the two-sided layers the loss is
$\tr(\Hout E \Hin E^\top)$, Eqs.~\eqref{eq:lq}--\eqref{eq:lk}. With
$\Hout^{-1} = U_{\mathrm{out}}^\top U_{\mathrm{out}}$ and rows quantised
in chunks, the exact OBS weight of a chunk's error matrix $E$ is
$\tr\big((U_{\mathrm{sub}}^\top U_{\mathrm{sub}})^{-1} E \Hin
E^\top\big)$; the solver already forms $U_{\mathrm{sub}}^{-1}$ for its row
update, so $r_i = \sum_k (U_{\mathrm{sub}}^{-1})_{ik}^2$ is free. It
reduces to $1/U_{\mathrm{out},ii}^2$ for chunk size 1. We use the
diagonal, i.e.\ we neglect the cross-row terms of the chunk, which is the
same approximation the per-element loss increment makes in every
one-sided solver.

\paragraph{Layer normalisation.} Gains are not comparable across layers:
they scale with the layer's Hessian magnitudes. We normalise by the
solver's own loss functional evaluated on the unquantised weight itself,
$N_\ell = \mathcal{L}_\ell(W_\ell)/\mathrm{numel}(W_\ell)$, computed once
per layer from the same damped Hessians the solver uses ($\sum_h
\tr(W_h\Hin W_h^\top)$ one-sided, $\sum_h \tr(\Hout W_h \Hin W_h^\top)$
two-sided), so that $\gain/N_\ell$ is dimensionless in every layer.

\paragraph{Selection rules (arms).}
\begin{itemize}\itemsep 0pt
\item \textbf{Arm A --- static Fisher (baseline).} $\mathrm{refine}_{ij}
  = M_{ij}$, $M$ the per-tensor top-$p$ mask by the SqueezeLLM diagonal
  Fisher $F_{ij} = \sum_n (\partial\mathcal{L}_n/\partial w_{ij})^2$ over
  the same calibration set. Needs one full backward pass over the model.
\item \textbf{Arm B --- in-loop, fixed per-layer $p$.} Top-$k$ gains per
  column with $k_j = \lfloor pn(j{+}1)\rfloor - \lfloor pnj\rfloor$ (so
  the mean is exactly $p$), requiring $\gain > 0$; over the rows present
  in the solver call (whole tensor, or one row chunk for two-sided
  layers).
\item \textbf{Arm C --- in-loop, one global $\lambda$.}
  $\mathrm{refine}_{ij} = [\gain_{ij}/N_\ell > \lambda]$ with one
  $\lambda$ for the whole model. This is the rule we advocate: it
  allocates the budget across layers.
\item \textbf{Arm D --- hybrid.} As C with $\gain_{ij} \leftarrow
  \gain_{ij} F_{ij}/\bar F_\ell$.
\item \textbf{Endpoints.} E3: never refine ($\lambda = \infty$); E4:
  always refine ($\lambda = 0$). Both run through the modified solver
  and bracket every arm.
\end{itemize}
The error propagated is that of the single value actually stored; we do
not propagate an average of the coarse and fine errors (MatGPTQ's rule,
which optimises a different, multi-target objective).

\subsection{Calibrating $\lambda$ on rate only}
\label{sec:lambda}

$\lambda$ is a rate knob and must never be tuned on test perplexity. Arm
B logs $\gain/N_\ell$ for every element at no extra cost; from the
model-wide empirical CDF we read off $\lambda$ at the quantile whose
expected refined fraction equals the target $p$. One arm-C run then
realises a rate; if it misses the target by more than 0.02 bpw we take a
secant step in $\log\lambda$ and rerun, at most twice. Because the arm-C
gains are computed on its own compensated weights, the realised fraction
differs slightly from the B-predicted one (\tbd{0.996\% vs.\ 1.000\% at
seed 0}); in our runs \tbd{zero} secant steps were needed. The same
$\lambda$ is reused across calibration seeds and the rate drift is
reported.

%%%%%%%%%%%%%%%%%%%%%%%%%%%%%%%%%%%%%%%%%%%%%%%%%%%%%%%%%%%%%%%%%%%%%%%%
\section{Experimental protocol}
\label{sec:protocol}

\paragraph{Model and solver.} Llama-3.2-1B \citep{grattafiori2024llama3}
(bf16, wiki2 PPL 9.751 at context 2048), quantised with TurboBoA
\citep{kim2025boa} at its released Llama defaults with $g{=}128$:
Hessian-weighted per-group scale search, attention-aware (two-sided)
solves for q/k with row chunks of 16, block-diagonal per-head value
Hessian, input-difference correction ($\alpha = 0.25$), 1\% damping, no
activation ordering (forced off by the solver for $g \neq -1$). Nothing
in this configuration is changed between arms.

\paragraph{Calibration and evaluation.} 128 sequences of 2048 tokens from
WikiText-2 train \citep{merity2017wikitext}, drawn by the solver's sampler
with seeds 0, 1, 2 (the same seed gives the same set to every arm). Test
perplexity on WikiText-2 test and C4-new \citep{raffel2020c4} at context
2048. The Fisher for arm A is accumulated in fp32 over the same 128
sequences of the same seed, on the transformer-block linears only.

\paragraph{Rate accounting and matching.} Per tensor one Huffman table
over the $2^{b+1}$ symbols ($2^b$ for native $b$-bit runs, which is the
fair rate for pure $b$-bit) plus table bits, 16-bit scales and $b$-bit
zeros per group; the script is solver-independent and takes only the
integer codes and group parameters. The rate target $R_1$ is the Huffman
bpw of arm A at $p = 1\%$, seed 0; arms B and C adjust $p/\lambda$ to
$R_1 \pm 0.02$, on rate only. A second point $R_2$ at $p \approx 4\%$
draws the rate--distortion line.

\paragraph{Decision rule (pre-registered).} With $\mu$ the mean over
three seeds and $\sigma$ the pooled seed-to-seed standard deviation of A,
B, C: PASS if $\mu_C - \mu_A \le -\max(2\sigma, 0.05)$ PPL; PARTIAL if
$|\mu_B - \mu_A| \le \sigma$ or $|\mu_C - \mu_A| \le \sigma$ (in-loop
matches Fisher at zero gradient cost); FAIL otherwise.

\paragraph{Correctness gates.} (a) E3 must be bit-identical to the native
$b$-bit run: codes equal $2\times$ the native codes, identical scales and
zeros, identical perplexity --- this proves the modified loop is a no-op
when nothing is refined. (b) E4 must be within a small margin of native
$(b{+}1)$-bit (the derived fine grid is not a native search). (c) Native
$b$-bit and $(b{+}1)$-bit rates must differ by $\approx 1$ bpw under the
accounting. (d) Arm B's perplexity must lie between E3 and E4, its
realised fraction must equal the target, and every refined element must
have positive gain.

\subsection{Experiments to run}

Table~\ref{tab:plan} (Appendix~\ref{app:plan}) lists every run the paper
needs, its purpose, and its status at the time of writing; the run tags
are the directory names under \texttt{results/runs/} in the artifact.



%%%%%%%%%%%%%%%%%%%%%%%%%%%%%%%%%%%%%%%%%%%%%%%%%%%%%%%%%%%%%%%%%%%%%%%%
\section{Results}
\label{sec:results}

\subsection{Endpoints and gates}
\label{sec:endpoints}

\tbd{Fill from results/SUMMARY.md ``Phase 0/1''.} Native W3:
\tbd{12.011 wiki2, 2.996 bpw (w3\_s0)}; native W4: \tbd{10.224, 3.902
(w4\_s0)}. E3 is bit-identical to W3 (\tbd{0 of 973,078,528 codes
differ, identical scales/zeros, identical PPL}) --- gate (a). E4:
\tbd{10.533 at 3.982 bpw}, i.e.\ \tbd{+0.31} PPL above native W4 --- gate
(b), the price of the derived fine grid. The rate accounting places W3
and W4 \tbd{0.906} bpw apart --- gate (c). Both native rates sit above the
2.6--2.9 / 3.3--3.6 ranges one might expect from peaked histograms:
TurboBoA's Hessian-weighted scale search uses the full code range, so its
histograms are broad and there is little entropy to remove from a pure
$b$-bit stream; the entropy saving in our scheme comes from the nesting,
not from the base quantizer.

\subsection{Main comparison at matched rate}
\label{sec:main}

Table~\ref{tab:main} is the main result. \tbd{One paragraph: $\mu_A$,
$\mu_B$, $\mu_C$, $\sigma$, the decision (PASS/PARTIAL/FAIL) under the
pre-registered rule, and the C4 numbers as a second metric. State the
rate drift of B and C across seeds at fixed $p/\lambda$.}

\begin{table*}[t]
\caption{Llama-3.2-1B, TurboBoA $g{=}128$, matched Huffman rate $R_1 \pm
0.02$. PPL at context 2048; mean $\pm$ std over calibration seeds 0/1/2
where three seeds exist. Nominal bpw includes the 0.15 bpw of
scales/zeros.}
\label{tab:main}
\vskip 0.1in
\centering
\small
\resizebox{\textwidth}{!}{%
\begin{tabular}{llcccccccc}
\toprule
Arm & selection & $p/\lambda$ & seeds & nominal & ideal bpw & Huffman bpw & wiki2 PPL & C4 PPL & extra pass \\
\midrule
FP16 & -- & -- & -- & 16 & -- & -- & 9.751 & 14.02 & -- \\
native W3 & -- & -- & 0 & 3.15 & \tbd{2.966} & \tbd{2.996} & \tbd{12.011} & \tbd{20.115} & -- \\
native W4 & -- & -- & 0 & 4.16 & \tbd{3.875} & \tbd{3.902} & \tbd{10.224} & \tbd{15.428} & -- \\
E3 & never & $\infty$ & 0 & 4.15 & \tbd{2.966} & \tbd{2.996} & \tbd{12.011} & \tbd{20.115} & -- \\
E4 & always & 0 & 0 & 4.15 & \tbd{3.946} & \tbd{3.982} & \tbd{10.533} & \tbd{16.346} & -- \\
A & Fisher top-1\% / tensor & 1\% & 0,1,2 & 4.15 & \tbd{3.008} & \tbd{3.071 $= R_1$} & \tbd{$\mu\pm\sigma$} & \blank & backward \\
B & in-loop top-$k$ / column & \tbd{1\%} & 0,1,2 & 4.15 & \tbd{3.027} & \tbd{3.076} & \tbd{$\mu\pm\sigma$} & \blank & none \\
C & in-loop global $\lambda$ & \tbd{0.0759} & 0,1,2 & 4.15 & \tbd{3.031} & \tbd{3.086} & \tbd{$\mu\pm\sigma$} & \blank & none \\
D & $(\gain\cdot F) > \lambda$ & \blank & 0,1,2 & 4.15 & \blank & \blank & \blank & \blank & backward \\
A-ref & existing pipeline (different solver) & 1\% & -- & -- & -- & \tbd{2.6} & \blank & -- & backward \\
\bottomrule
\end{tabular}}
\end{table*}

\paragraph{Odd codes vs.\ refined elements.} Arm A refines exactly
1.000\% of every tensor but only \tbd{0.46\%} of its codes end up odd:
\tbd{54\%} of the Fisher picks are free elements
(Section~\ref{sec:free}). Arms B and C, which require $\gain > 0$,
produce \tbd{$\approx 2\times$} the odd codes at the same refined
fraction. This is why we match on rate: at matched Huffman rate the
in-loop arms store \tbd{twice} as many effective refinements as the
Fisher arm for the same bits. (A caveat we report: integer Huffman
lengths do not resolve a 0.5\% change in the odd-symbol mass, so B/C's
Shannon rate is \tbd{0.02} bpw above A's while their Huffman rates agree
to \tbd{0.005}; an ANS coder would make the comparison slightly less
favourable to B/C by that amount.)

\begin{figure}[t]
\centering
\includegraphics[width=\columnwidth]{figs/fig1_refined_fraction.png}
\caption{Refined fraction per layer, seed 0, arms A/B/C at matched rate.
\tbd{Regenerate with C included: \texttt{scripts/figs.py --A A\_p1\_s0
--B B\_p1\_s0 --C C\_l0.0759\_s0}.}}
\label{fig:alloc}
\end{figure}

\begin{figure}[t]
\centering
\includegraphics[width=\columnwidth]{figs/fig2_jaccard.png}
\caption{Jaccard(A-mask, B-mask) per layer, seed 0.}
\label{fig:jaccard}
\end{figure}

\paragraph{Where the budget goes.} Figure~\ref{fig:alloc} shows the
refined fraction per layer. Arm B is flat by construction. Arm C, with
one global $\lambda$, gives the query/key projections \tbd{$\approx
0.005\%$} and spends the budget on the output and down projections
(\tbd{up to 3.9\%} in the middle blocks), yet ends \tbd{lowest} in
perplexity: the two-sided loss weight reports that fine-grid errors in
q/k are cheap relative to those layers' own loss scale. \tbd{Confirm with
the $N_\ell = 1$ ablation (4c) whether the normalisation is what makes
this allocation work.}

\paragraph{The two rules pick different weights.} The per-layer Jaccard
index between the Fisher mask and the in-loop mask is \tbd{0.009 mean,
0.047 max} (Figure~\ref{fig:jaccard}) --- chance level for two 1\% masks
is $\approx 0.005$. The two criteria are nearly independent, and both
help; the in-loop one helps more per stored bit.

\paragraph{Rate--distortion.} Figure~\ref{fig:rd} places every run on the
(Huffman bpw, PPL) plane. \tbd{Add the $R_2$ point (4b) and the 2$\to$3-bit
runs; the line from E3 through A/B/C at $R_1$ and $R_2$ toward E4 is the
deliverable.}

\begin{figure}[t]
\centering
\includegraphics[width=0.9\columnwidth]{figs/fig3_rate_distortion.png}
\caption{Rate--distortion, Llama-3.2-1B. \tbd{Regenerate with C, $R_2$
and the 2-bit runs.}}
\label{fig:rd}
\end{figure}

\subsection{Ablations}

\paragraph{Fisher on top of the gain (arm D).} \tbd{Table row + one
sentence: does multiplying the in-loop gain by the Fisher change the
outcome? If D $\approx$ C, the gradient pass buys nothing.}

\paragraph{Layer normalisation ($N_\ell = 1$).} \tbd{One sentence: rate
and PPL of arm C with raw gains and the same target rate; per-layer
allocation plot.}

\paragraph{Second rate point $R_2$.} \tbd{A/B/C at $p \approx 4\%$: does
the gap persist / widen with budget?}

\subsection{One level down: 2$\to$3 bits}

The construction is bit-width agnostic; we repeat phases 0--3 with a
2-bit coarse grid and a 3-bit fine grid (codes 0..7, $s_f = s_c/2$, $z_f
= 2z_c$), seed 0, $p = 1\%$. \tbd{One paragraph. Expectation: the 2-bit
base is far lossier, so 1\% refinement moves PPL less in relative terms
and the free fraction may differ; report whether the ordering A/B/C is
preserved.}

\begin{table}[t]
\caption{2$\to$3-bit nesting, Llama-3.2-1B, seed 0.}
\label{tab:2to3}
\vskip 0.1in
\centering
\footnotesize
\begin{tabular}{lccc}
\toprule
Arm & Huff.\ bpw & wiki2 & C4 \\
\midrule
native W2 & \blank & \blank & \blank \\
E2 (coarse only) & \blank & \blank & \blank \\
E3 (all fine, derived) & \blank & \blank & \blank \\
native W3 & \tbd{2.996} & \tbd{12.011} & \tbd{20.115} \\
A (Fisher 1\%) & \blank & \blank & \blank \\
B (in-loop top-$k$) & \blank & \blank & \blank \\
C (in-loop global $\lambda$) & \blank & \blank & \blank \\
\bottomrule
\end{tabular}
\end{table}

\subsection{Cost}
\label{sec:cost}

The in-loop rule adds \tbd{$\approx 150$ s} of per-column Python
bookkeeping to a \tbd{1060 s} solve on one A100 (all of it logging; the
selection itself is a top-$k$ or a compare on a vector of $\le 2048$
elements per column) and removes the Fisher pass (\tbd{25 s} on an A100
for a 1B model, but a full backward pass over the model with a 12.7 GB
peak, which for larger models is the step that no longer fits next to
the solver).

% ---- new multi-model sections (6.6 - 6.11) ----
\input{sec6_extended}

%%%%%%%%%%%%%%%%%%%%%%%%%%%%%%%%%%%%%%%%%%%%%%%%%%%%%%%%%%%%%%%%%%%%%%%%
\section{Bitstream format and encoder}
\label{sec:bitstream}

\teammate{Write this section. It must specify: (i) the on-disk/in-memory
layout of a quantised tensor: per-tensor Huffman table ($2^{b+1}$ code
lengths, canonical Huffman), per-group 16-bit scale and $b$-bit zero, the
coded stream and its block/chunk structure for random access (e.g.\ one
independently decodable chunk per $N$ rows or per tile so that a GEMM
tile never spans two chunks); (ii) how the code stream is produced from
the solver's integer codes (\texttt{results/runs/<tag>/layers.*.npz}:
codes uint8 $[d_{\mathrm{out}}, d_{\mathrm{in}}]$, scale fp32
$[d_{\mathrm{out}}, n_{\mathrm{groups}}]$, zero uint8, bits, zero\_bits);
(iii) the exact bits-per-weight of the format including all headers, and
how close it is to the Huffman number reported in Table~\ref{tab:main}
(the paper's rate axis); (iv) whether an ANS/rANS variant is used or
considered and its rate vs.\ decode-speed trade-off; (v) the choice of
the fine/coarse zero and scale precision.}

%%%%%%%%%%%%%%%%%%%%%%%%%%%%%%%%%%%%%%%%%%%%%%%%%%%%%%%%%%%%%%%%%%%%%%%%
\section{Decoder and fused dequantisation kernels}
\label{sec:kernels}

\teammate{Write this section. It must cover: (i) the Huffman decode
strategy on GPU (table-driven canonical decode, bits consumed per step,
warp-level parallelism, how chunk boundaries give independent decode
streams); (ii) the fused decode $\to$ dequantise ($\hat w = s_f(q -
z_f)$, no branch on refinement) $\to$ GEMV/GEMM kernel for the
memory-bound decode phase, and how it compares to a 4-bit packed kernel
(e.g.\ Marlin \citep{frantar2025marlin}, LUT-GEMM
\citep{park2024lutgemm}) that reads 4 bits/weight; (iii) the prefill path
(decode once to a 4-bit packed buffer vs.\ decode on the fly); (iv)
integration into the serving stack used for the measurements; (v) any
restrictions the kernel imposes on the format (alignment, chunk size)
that had to be fed back into Section~\ref{sec:bitstream}.}

%%%%%%%%%%%%%%%%%%%%%%%%%%%%%%%%%%%%%%%%%%%%%%%%%%%%%%%%%%%%%%%%%%%%%%%%
\section{Systems evaluation}
\label{sec:sysval}

\subsection{Llama-3.2-1B on one GPU}

\teammate{Write this subsection. Required measurements: (i) kernel
micro-benchmarks --- decode throughput (weights/s and GB/s of compressed
bytes read) vs.\ a 4-bit packed baseline and vs.\ fp16, for the layer
shapes of Llama-3.2-1B; (ii) end-to-end tokens/s at batch 1 and small
batches, fp16 vs.\ 4-bit vs.\ this format, on \tbd{GPU(s)}; (iii) model
memory footprint on device; (iv) the compression/encode time; (v) a
breakdown showing whether Huffman decode is hidden under the memory-bound
GEMV or adds latency, and at what batch size it stops being hidden.
Include a table and one figure.}

% ---- new systems-at-scale subsection (9.2) ----
\input{sec9_extended}

%%%%%%%%%%%%%%%%%%%%%%%%%%%%%%%%%%%%%%%%%%%%%%%%%%%%%%%%%%%%%%%%%%%%%%%%
\section{Related work}
\label{sec:related}

\paragraph{Error-compensating PTQ.} OBS \citep{hassibi1993obs}, OBC
\citep{frantar2022obc} and GPTQ \citep{frantar2023gptq} are the lineage of
the column loop we modify; AWQ \citep{lin2024awq} and OmniQuant
\citep{shao2024omniquant} act on scales rather than rounding, and GPTAQ
\citep{li2025gptaq} corrects the calibration target. BoA/TurboBoA
\citep{kim2025boa} adds an output-side Hessian for attention projections;
our row weight $r_i$ comes from it.

\paragraph{Element- and group-level mixed precision.} SpQR
\citep{dettmers2024spqr} detects outliers in-loop from the OBS error but
stores them as an explicit FP16 sparse structure; our score is the
difference $e_c^2 - e_f^2$, not $e_c^2$, and the refined values live in
the same code stream. SqueezeLLM \citep{kim2024squeezellm} and OWQ
\citep{lee2024owq} select by static sensitivity outside the solver (our
arm A); SliM-LLM \citep{huang2024slimllm} assigns bit-widths per group by
salience; GuidedQuant \citep{kim2025guidedquant} weights layer-wise
objectives by end-loss gradients. LLM.int8() \citep{dettmers2022llmint8}
separates outlier columns.

\paragraph{Nested and Matryoshka codes.} NestQuant
\citep{savkin2025nestquant} uses a fixed nested $E_8$ pair per tensor
with no per-element choice. Matryoshka Quantization
\citep{nair2025matquant}, Any-Precision LLM \citep{park2024anyprecision}
and MatGPTQ \citep{kleinegger2026matgptq} store one master code per
element and serve lower precisions by MSB slicing; MatGPTQ chooses the
master code by brute force over the weighted sum of errors at all target
widths, propagates the average error, and varies precision per layer via
EvoPress \citep{sieberling2024evopress}. None chooses a width per
element, keeps a per-element mask, or uses the nesting for entropy coding;
MatGPTQ's MSB slice $\mathrm{round}(q/2^{c-r})\,2^{c-r}$ is the same
sub-lattice as our even codes, so MatGPTQ-EP at $\approx 3.0$--$3.25$ bpw
is the natural baseline row (Table~\ref{tab:published}).

\paragraph{Vector and lattice quantization.} QuIP\#
\citep{tseng2024quipsharp}, QTIP \citep{tseng2024qtip}, AQLM
\citep{egiazarian2024aqlm} and GPTVQ \citep{vanbaalen2024gptvq} reach
lower rates with codebooks and incoherence processing, at the cost of a
lookup-based kernel; our scheme is orthogonal (it changes only which
scalar level is stored) and keeps a multiply--add dequantiser.

\paragraph{Entropy coding of weights.} Deep Compression
\citep{han2016deepcompression} Huffman-coded codebook indices after
clustering; DFloat11 \citep{zhang2025dfloat11}, NeuZip
\citep{hao2024neuzip} and ZipNN \citep{hershcovitch2024zipnn} entropy-code
floating-point exponents losslessly, and ZipServ \citep{zipserv2026} and
\citet{shannonbound2026} fuse the decode into the GEMM. We entropy-code a
stream whose low entropy is designed in by the nesting.

%%%%%%%%%%%%%%%%%%%%%%%%%%%%%%%%%%%%%%%%%%%%%%%%%%%%%%%%%%%%%%%%%%%%%%%%
\section{Conclusion}
\label{sec:conclusion}

A 3-bit grid nested in a 4-bit grid lets a PTQ solver give individual
weights one extra bit with no mask, no index and no second scale, and
lets the solver itself decide which weights get it, from the error weight
it already computes on the residual it already holds. At matched
entropy-coded rate the in-loop rule \tbd{matches/beats} a static Fisher
mask on Llama-3.2-1B while removing the gradient pass, and one global
threshold allocates the budget across layers. \tbd{One sentence on the
2-bit result.} \tbd{One sentence on 3B/8B and the other families.}
\teammate{One sentence on the systems result.} Limitations: we evaluated
one solver family within a fixed compute budget; the derived fine grid
costs \tbd{0.3} PPL relative to a native search when all weights are
refined, which bounds the benefit at large budgets; Huffman rate matching
is coarse at the 0.01 bpw level; and \tbd{the 70B point is not budgeted
/ was run on arms E3/A/C only}.

%%%%%%%%%%%%%%%%%%%%%%%%%%%%%%%%%%%%%%%%%%%%%%%%%%%%%%%%%%%%%%%%%%%%%%%%
\bibliography{refs}
\bibliographystyle{mlsys2025}

%%%%%%%%%%%%%%%%%%%%%%%%%%%%%%%%%%%%%%%%%%%%%%%%%%%%%%%%%%%%%%%%%%%%%%%%
% MLSys: the appendix is uploaded as a separate file. Keep it here for
% drafting; split at submission time.
\newpage
\appendix
\onecolumn

\section{Experiment plan}
\label{app:plan}

\begin{table}[h]
\caption{Experiment plan. Seed = calibration seed. Status: $\checkmark$
done, $\circ$ running/queued, -- not started. Run tags are directory
names under \texttt{results/runs/} in the artifact.}
\label{tab:plan}
\vskip 0.1in
\centering
\small
\resizebox{\textwidth}{!}{%
\begin{tabular}{clllc}
\toprule
Phase & Run(s) & Purpose & Seeds & Status \\
\midrule
0 & w3\_s0, w4\_s0 & native 3/4-bit reference, wall-clock & 0 & $\checkmark$ \\
1 & e3\_s0, e4\_s0 & endpoints; gates (a)--(c) & 0 & $\checkmark$ \\
2 & Fisher; A\_p1\_s0; B\_p1\_s0 & $R_1$; in-loop top-$k$ at $R_1$; Jaccard(A,B); gain histograms & 0 & $\checkmark$ \\
3 & $\lambda$ from B's CDF; C\_l*\_s0; A/B/C seeds 1, 2 & main table; pre-registered decision rule & 0,1,2 & \tbd{$\circ$ seed 2 pending} \\
4a & arm D at $R_1$ & does Fisher add anything on top of the in-loop gain? & 0,1,2 & -- \\
4b & A/B/C at $R_2$ ($p \approx 4\%$) & two-point rate--distortion line & 0 & -- \\
4c & arm C with $N_\ell = 1$ & is layer normalisation what makes one global $\lambda$ work? & 0 & -- \\
5 & 2$\to$3-bit mirror of phases 0--3 (w2, E2/E3 endpoints, A, B, C) & does the rule transfer one level down? & 0 & -- \\
6 & $g = -1$ (per-channel) point for E3/A/B/C & robustness to the grouping choice (needs act-order support for the mask) & 0 & to decide \\
7 & Llama-3.2-3B / Llama-3.1-8B, phases 0--3 & scaling; recovery ratio (Section~\ref{sec:scaling}) & 0 & -- \\
8 & zero-shot tasks (lm-eval) for the final A/B/C checkpoints & downstream sanity (Section~\ref{sec:zeroshot}) & 0 & -- \\
9 & OPT-6.7B, Mistral-7B-v0.3, Qwen3-8B: E3/A/C & other families; BoA / MatGPTQ cross-check (Section~\ref{sec:families}) & 0 & -- \\
10 & Llama-2-7B/13B full pipeline; baseline re-runs (GPTQ, AWQ, SliM-LLM, GPTAQ); rate script on their codes & published comparison (Section~\ref{sec:published}) & 0 & -- \\
11 & sensitivity blocks ($p$ sweep, $g$, calibration, solver, $r_i$-in-gain, A$'$, coder) & reviewer questions (Section~\ref{sec:sensitivity}) & 0 & -- \\
12 & cost table across sizes & scaling argument (Section~\ref{sec:cost-scale}) & 0 & -- \\
S & \teammate{kernel microbenchmarks, end-to-end decode throughput, memory} & systems half & -- & -- \\
\bottomrule
\end{tabular}}
\end{table}

\section{Reproducibility}
\label{app:repro}

\tbd{Artifact: repository layout (\texttt{turboboa/} vendored solver with
the flag-guarded patch, \texttt{scripts/} one script per step,
\texttt{results/} per-run JSON + SUMMARY.md), exact commands from
HANDOFF.md, hardware (1$\times$ A100-PCIE-40GB, torch 2.6.0+cu124,
transformers 4.53.0), wall-clock per run ($\approx$ 20 min quant + 3 min
eval).}

\section{Additional figures}
\label{app:figs}

\tbd{Gain histograms per layer (log $x$), free fraction per layer,
per-layer table (refined fraction, ideal entropy, Huffman bpw, $N_\ell$,
wall time) for the final A/B/C runs. Per-layer allocation of arm C at 3B
and 8B (Figure~\ref{fig:alloc-scale} data). Full lm-eval per-task tables
behind Table~\ref{tab:zeroshot}.}

\input{app_sensitivity}

\section{Rate accounting details}
\label{app:rate}

\tbd{Worked example of Eq.~\eqref{eq:rate} on one tensor: histogram,
Huffman lengths, table bits, scale/zero bits; the same for a native
3-bit tensor with an 8-symbol table; and the tANS rate for the same
stream.}

\end{document}
